
\documentclass{article}
\usepackage{amsmath} % Pour les environnements d'équations

\title{Équations de Maxwell}
\author{Agent de Coding}
\date{\today}

\begin{document}

\maketitle

\section*{Introduction}
Les équations de Maxwell sont un ensemble de quatre équations fondamentales qui décrivent le comportement des champs électrique et magnétique, ainsi que leurs interactions avec la matière. Elles sont la base de l'électromagnétisme classique.

\section*{Une Équation de Maxwell : Équation de Gauss pour l'Électrostatique}

L'équation de Gauss pour l'électrostatique est l'une des quatre équations de Maxwell. Elle relie le flux du champ électrique à travers une surface fermée à la charge électrique totale contenue dans cette surface.

En forme intégrale, elle s'écrit :
\begin{equation*}
    \oint_S \mathbf{E} \cdot d\mathbf{A} = \frac{Q_{\text{int}}}{\varepsilon_0}
\end{equation*}
où :
\begin{itemize}
    \item $\mathbf{E}$ est le champ électrique.
    \item $d\mathbf{A}$ est un vecteur élémentaire de surface, dont la direction est normale à la surface $S$.
    \item $Q_{\text{int}}$ est la charge électrique totale contenue à l'intérieur de la surface fermée $S$.
    \item $\varepsilon_0$ est la permittivité du vide.
\end{itemize}

En forme différentielle, elle s'écrit :
\begin{equation*}
    \nabla \cdot \mathbf{E} = \frac{\rho}{\varepsilon_0}
\end{equation*}
où :
\begin{itemize}
    \item $\nabla \cdot \mathbf{E}$ est la divergence du champ électrique.
    \item $\rho$ est la densité de charge électrique volumique.
\end{itemize}

Cette équation exprime que les charges électriques sont la source des champs électriques.

\end{document}
